%%%%%%%%%%%%%%%%%%%%%%%%%%%%%%%%%%%%%%%%%%%%%%%%%%%%%%%%%%%%%%%%%%%%%%%%%%%
%%  Project: CPHOS LaTeX Template
%%  Copyright © 2026 gameswu & CPHOS (www.cphos.cn)
%%
%%  Permission is hereby granted, free of charge, to any person obtaining a 
%%  copy of this software and associated documentation files (the "Software"), 
%%  to deal in the Software without restriction, including without limitation 
%%  the rights to use, copy, modify, merge, publish, distribute, sublicense, 
%%  and/or sell copies of the Software, and to permit persons to whom the 
%%  Software is furnished to do so, subject to the following conditions:
%%
%%  The above copyright notice and this permission notice shall be included 
%%  in all copies or substantial portions of the Software.
%%
%%  THE SOFTWARE IS PROVIDED "AS IS", WITHOUT WARRANTY OF ANY KIND, EXPRESS 
%%  OR IMPLIED, INCLUDING BUT NOT LIMITED TO THE WARRANTIES OF MERCHANTABILITY, 
%%  FITNESS FOR A PARTICULAR PURPOSE AND NONINFRINGEMENT. IN NO EVENT SHALL 
%%  THE AUTHORS OR COPYRIGHT HOLDERS BE LIABLE FOR ANY CLAIM, DAMAGES OR OTHER 
%%  LIABILITY, WHETHER IN AN ACTION OF CONTRACT, TORT OR OTHERWISE, ARISING 
%%  FROM, OUT OF OR IN CONNECTION WITH THE SOFTWARE OR THE USE OR OTHER 
%%  DEALINGS IN THE SOFTWARE.
%%%%%%%%%%%%%%%%%%%%%%%%%%%%%%%%%%%%%%%%%%%%%%%%%%%%%%%%%%%%%%%%%%%%%%%%%%%

% CPHOS理论命题模板

\documentclass[answer]{cphos}

% === 设置文档信息 ===
\cphostitle{CPHOS物理竞赛联考}
\cphossubtitle{理论试题命题模板}

\setscorecheck{true}  % 启用分值检查，可设为false关闭

\begin{document}
\begin{problem}[40]{激光制冷技术}

% === 题干 ===
\begin{problemstatement}
激光制冷（Laser Cooling）是一种利用激光与原子或分子相互作用来降低物质温度的技术，广泛应用于原子物理、量子信息等领域。通过精确调控激光频率和强度，可以实现对原子或分子的有效冷却，甚至达到接近绝对零度的温度。

设原子某一能级的跃迁频率为$\nu$，当一束轻微红失谐的激光照射在原子上时（激光频率应略小于$\nu$，本题中可将其近似为$\nu$），原子更倾向于吸收来自与其运动方向相反的激光光子，从而减慢其运动速度。原子吸收一个光子后会立即发生自发辐射，在这个过程中，原子会随机且均匀地向任意方向发射一个频率为$\nu$的光子，导致平均意义上发射过程中的动量变化为零。因此，原子在吸收光子时动量会减小，而在自发辐射时动量变化平均为零，实现冷却效果。

当原子在激光射来的方向上的速度投影为$u$时，该吸收-发射循环的频率可以表示为
\begin{equation}
    R(u) = \varGamma(1+\kappa u) \eqtag{1}
\end{equation}
其中$\varGamma$是原子在该能级的自然衰变率，$\kappa$是一个与激光频率$\nu$和原子质量$m$相关的常数。当有多束激光从不同方向照射时，各激光均会对应一个由该方向速率决定的吸收-发射循环频率。

\vspace{1em}
\noindent\textbf{注意：原子激光制冷技术中，激光通常在近红外/可见光/近紫外波段，光子能量与原子相对论静能的量级关系为}
$$
\frac{h\nu}{m_pc^2}=\frac{h}{m_pc\lambda}\sim 10^{-9}
$$
\noindent\textbf{因此本题可以忽略相对论效应的影响。}

\vspace{1em}
\subq{1} 本问中只考虑原子在$z$方向上的运动，有两束制冷激光分别从$+z$和$-z$方向照射原子。

\subsubq{1.1} 将原子吸收光子的过程看作近似连续的，证明原子的受到的阻尼力与速度的关系满足
\begin{equation}
    F_z = -\gamma mu_z \eqtag{2}
\end{equation}
并使用$\nu,\varGamma,\kappa,m$以及光速$c$、普朗克常量$h$来表示阻尼系数$\gamma$。

\subsubq{1.2} 在真实情况下，原子吸收光子的过程是离散的。原子的初始速度为$u_z$，当原子吸收了一个从$+z$方向来的光子后，计算原子的动能改变量$\Delta E$（不考虑接下来的发射过程影响）。

\subq{2} 本问中考虑原子在三维空间中的运动，六束制冷激光分别从$\pm x,\pm y,\pm z$方向照射原子。

\subsubq{2.1} 暂不考虑发射光子产生的影响，设原子的方均速率为$\langle u^2 \rangle$，表示出原子在六束激光制冷下单位时间的动能变化率$\langle \dot{E}_c \rangle$。

\subsubq{2.2} 接着考虑发射光子产生的影响，原子随机且均匀地向任意方向发射频率为$\nu$的光子，该过程可以视为动量空间中的随机游走。设原子在动量空间进行了$N$步随机游走：$\vec{p}_1, \vec{p}_2, \cdots, \vec{p}_N$，由于是均匀随机的，累积的动量改变为零，但是动量改变的方均值不为零且满足
\begin{equation}
    \Delta p = \vec{p}_1 + \vec{p}_2 + \cdots + \vec{p}_N = 0 \eqtag{3}
\end{equation}
\begin{equation}
    (\Delta p)^2 = p_1^2 + p_2^2 + \cdots + p_N^2 \eqtag{4}
\end{equation}
试求原子在发射光子产生的随机游走影响下单位时间的动能增加率$\langle \dot{E}_h \rangle$。

\subsubq{2.3} 结合前两小问的结果，原子的动能总变化率为$\langle \dot{E} \rangle = \langle \dot{E}_c \rangle + \langle \dot{E}_h \rangle$，试推导出激光制冷能达到的最低极限温度$T_{\text{limit}}$。

\end{problemstatement}

% === 解答 ===
\begin{solution}
\solsubq{1}{11}
\solsubsubq{1.1}{6}
在单位时间内，原子从$+z$方向吸收光子的次数为$R(u_z)$，从$-z$方向吸收光子的次数为$R(-u_z)$。每次吸收一个光子，原子动量改变$\Delta p = h\nu/c$，发射光子的平均动量改变为零，因此原子在z方向受到的力为
\begin{equation}
    F_z =\frac{\mathrm{d}p}{\mathrm{d}t}\approx \frac{\Delta p}{\Delta t}=-\frac{h\nu}{c} [R(u_z) - R(-u_z)] \eqtagscore{1}{3}
\end{equation}
将$R$代入上式，得到
\begin{equation}
    F_z = -\frac{h\nu}{c} \varGamma [\kappa u_z - (-\kappa u_z)] = -\frac{2h\nu\varGamma\kappa}{c} u_z \eqtag{2}
\end{equation}
故阻尼系数为
\begin{equation}
    \gamma = \frac{2h\nu\varGamma\kappa}{mc} \eqtagscore{3}{3}
\end{equation}

\solsubsubq{1.2}{5}
记吸收光子造成的动量改变量为$\delta p_z = -h\nu/c$，则原子的动能改变量为
\begin{equation}
    \Delta E = \frac{1}{2}m\left( u_z+\frac{\delta p_z}{m} \right)^2-\frac{1}{2}mu_z^2=-\frac{u_zh\nu}{c}+\frac{(h\nu)^2}{2mc^2} \eqtagscore{4}{5} \label{eq:4}
\end{equation}

\solsubq{2}{29}
\solsubsubq{2.1}{9}
由\ref{eq:4}可知，$+z$方向射来的激光造成的动能变化率为
\begin{equation}
    \dot{E}_{+z} = R(u_z) \Delta E = \varGamma(1+\kappa u_z) \left(-\frac{u_zh\nu}{c}+\frac{(h\nu)^2}{2mc^2}\right) \eqtagscore{5}{4}
\end{equation}
类似地可以写出其他五束激光造成的动能变化率，最终得到原子在六束激光制冷下单位时间的动能变化率为
\begin{equation}
    \begin{aligned}
    \langle \dot{E}_c \rangle &= \langle \dot{E}_{+x} \rangle + \langle \dot{E}_{-x} \rangle + \langle \dot{E}_{+y} \rangle + \langle \dot{E}_{-y} \rangle + \langle \dot{E}_{+z} \rangle + \langle \dot{E}_{-z} \rangle \\
    &= -\frac{2h\nu\varGamma\kappa}{c} \langle u^2 \rangle + \frac{3(h\nu)^2\varGamma}{mc^2}
    \end{aligned}
    \eqtagscore{6}{5}
\end{equation}

\solsubsubq{2.2}{9} 由题意，原子在动量空间中的随机游走满足$\Delta p = 0$和$(\Delta p)^2 = N(\delta p)^2$，其中$\delta p = h\nu/c$每一次发射光子造成的动量改变量。原子每吸收一个光子就会发射一个光子，因此单位时间内的发射次数为
\begin{equation}
    R_{tot} = R(u_x) + R(-u_x) + R(u_y) + R(-u_y) + R(u_z) + R(-u_z) = 6\varGamma \eqtagscore{7}{4}
\end{equation}
因此单位时间内的动量空间随机游走步数为$N = R_\text{tot}$，原子在发射光子产生的随机游走影响下单位时间的动能增加率为
\begin{equation}
    \langle \dot{E}_h \rangle = \frac{(\Delta p)^2}{2m} = \frac{N(\delta p)^2}{2m} = \frac{3(h\nu)^2\varGamma}{mc^2} \eqtagscore{8}{5}
\end{equation}

\solsubsubq{2.3}{11} 结合前两小问的结果，原子的动能总变化率为
\begin{equation}
    \langle \dot{E} \rangle = \langle \dot{E}_c \rangle + \langle \dot{E}_h \rangle = -\frac{2h\nu\varGamma\kappa}{c} \langle u^2 \rangle + \frac{6(h\nu)^2\varGamma}{mc^2} \eqtagscore{9}{3}
\end{equation}
当$\langle \dot{E} \rangle = 0$时，原子达到最低极限温度，此时方均速率为
\begin{equation}
    \langle u^2 \rangle = \frac{3h\nu}{mc\kappa} \eqtagscore{10}{4}
\end{equation}
根据动能与温度的关系$E = \frac{1}{2}m\langle u^2 \rangle = \frac{3}{2}k_B T$，可以得到激光制冷能达到的最低极限温度为
\begin{equation}
    T_{\text{limit}} = \frac{h\nu}{k_B c \kappa} \eqtagscore{11}{4}
\end{equation}

\scoring  % 评分标准自动使用题目环境中声明的总分

\end{solution}

\end{problem}

\end{document}
